%% ma: L12-B01-0037
%% lop: 12
%% chuong: 1
%% bai: 1
%% dang: 12.1.23
%% loai: TN4
%% muc: TH
%% dap_an: "D"
%% nguon: "HK1-DeOn_1_1.pdf (M10, Đề ôn tập học kì 1 Toán 12, Đề số 2), Phần I Câu 12"
%% kien_thuc: "Bài 1 (tính đơn điệu và cực trị)"
%% trang_thai: cho-duyet
%% ly_do: "Dạng mới đề xuất 12.1.23"
\begin{ex}
Đồ thị hàm số nào trong bốn hàm số liệt kê ở các phương án A, B, C, D dưới đây có đúng hai điểm cực trị?
\choice{$y=\dfrac{x-2}{-x+1}$}{$y=-x^3-x+3$}{$y=\dfrac{x^2-x-3}{x-1}$}{\True $y=-x^3+4x+3$}
\loigiai{A: hàm phân thức bậc nhất trên bậc nhất không có cực trị. B: $y'=-3x^2-1<0$. C: $y'=\dfrac{x^2-2x+4}{(x-1)^2}>0$. D: $y'=-3x^2+4$ có hai nghiệm phân biệt và đổi dấu nên có đúng hai điểm cực trị.}
\end{ex}
